\subsection{渐近展开的复合}

设 f 是 \(\mathcal{H}(\mathfrak{F},\mathbf{R})\) （相应地， \(\mathcal{H}(\mathfrak{F},\mathbf{C})\) ）中的一个函数，它容许一个相对于尺度 E 的、精确到 \(g_{\alpha}\) 的渐近展开，并且沿以 F 为基的滤子有极限 0。另一方面，设 h 是一个取值于 R（相应地， C）上的赋范空间 V 中的函数，定义在 R（相应地， C）中点 0 的一个邻域上，且在这个邻域上 n 次可导；则

\[
\mathbf {h} (t) = \mathbf {c} _ {0} + \mathbf {c} _ {1} t + \dots + \mathbf {c} _ {n} t ^ {n} + o (t ^ {n})
\]

在这个邻域上成立 (I, p.~21)，由此，在 \(\mathfrak{F}\) 的一个适当的集上

\[
\mathbf {h} \circ f = \mathbf {c} _ {0} + \mathbf {c} _ {1} f + \dots + \mathbf {c} _ {n} f ^ {n} + o (f ^ {n}).
\]

我们已在第 3 节中看到，如何为 \(c_{0} + c_{1}f + \cdots + c_{n}f^{n}\) 构造一个精确到 \(g_{\rho}\) 的渐近展开，其精确度由 f 的展开的精确度确定；此外，设 f 的渐近展开的系数不全为零，且 \(a_{\gamma}g_{\gamma}\) 是 f 的主部，并设 \(g_{\sigma} = g_{\gamma}^{n}\) ；若 \(\sigma < \rho\) ，则将得到 \(h \circ f\) 的一个精确到 \(g_{\sigma}\) 的展开（通过把 \(\sum_{k=0}^{n} c_{k}f^{k}\) 的展开限制到这个精确度）；反之，若 \(\rho \leqslant \sigma\) ，则 \(\sum_{k=0}^{n} c_{k}f^{k}\) 的展开也是 \(h \circ f\) 的一个精确到 \(g_{\rho}\) 的展开。

若 \(f\) 的渐近展开的所有项为零，且 \(g_{\alpha} \ll 1\) ，则 \(f \ll g_{\alpha}\) ，从而 \(f^k \ll g_{\alpha}^k \ll g_{\alpha}\) （对每个整数 \(k > 0\) ）；若 \(\mathbf{c}_m\) 是指标 \(> 0\) 的第一个不为零的系数（假设 \(\mathbf{c}_k\) 对指标 \(k > 0\) 不全为 0），则 \(\mathbf{c}_0\) 是 \(\mathbf{h} \circ f\) 的一个精确到 \(g_{\alpha}^{m}\) 的渐近展开。

在本节的其余部分，我们将限于 \(\mathcal{E}\) 中的函数取实值且在 \(\mathfrak{F}\) 的某个集上严格正的情形，并且只考虑 \(\mathcal{H}(\mathfrak{F},\mathbf{R})\) 中函数的渐近展开。首先假设对每个函数 \(g\in \mathcal{E}\) 和每个实数 \(\nu\) ， \(g^{\nu}\) 仍属于 \(\mathcal{E}\) ：这个条件例如被 \(x^{\alpha}\) 的尺度或 \(x^{\alpha}|\log x|^{\beta}\) 的尺度（ \(\alpha\) 与 \(\beta\) 为任意实数）在 \(+\infty\) 的一个邻域上或在 \(\mathbf{R}\) 中 0 的一个邻域上满足。这个性质蕴含 \(\mathcal{E}\) 中两个函数的商仍属于 \(\mathcal{E}\) 。这样，从相对于 \(\mathcal{E}\) 的、函数 \(f\in \mathcal{H}(\mathfrak{F},\mathbf{R})\) 的一个精确到 \(g_{\alpha}\) 的渐近展开，可导出 \(|f|^{\nu}\) 的一个展开（对每个实数 \(\nu\) ）。让我们限于 \(f\) 的展开的系数不全为零的情形，并设 \(a_{\gamma}g_{\gamma}\) 是 \(f\) 的主部；可以写 \(|f|^{\nu} = |a_{\gamma}|^{\nu}g_{\gamma}^{\nu}(1 + h)^{\nu}\) ，其中

\[
h = \sum_ {\gamma <   \lambda \leqslant \alpha} \frac {a _ {\lambda}}{a _ {\gamma}} \frac {g _ {\lambda}}{g _ {\gamma}} + o \left(\frac {g _ {\alpha}}{g _ {\gamma}}\right).
\]

在我们的假设下， \(\sum_{\gamma < \lambda \leqslant \alpha} \frac{a_{\lambda}}{a_{\gamma}} \frac{g_{\lambda}}{g_{\gamma}}\) 是 \(h\) 的一个渐近展开（精确到 \(g_{\alpha} / g_{\gamma}\) ）；由于 \(h\) 沿 \(\mathfrak{F}\) 趋于 0，上面描述的方法给出 \((1 + h)^{\nu}\) 的一个渐近展开，然后得到 \(|f|^{\nu}\) 的一个展开（通过乘以 \(\left|a_{\gamma}\right|^{\nu} g_{\gamma}^{\nu}\) ）。

对 \(f\) 用同样的假设，可以写

\[
\log | f | = \log \left| a _ {\gamma} g _ {\gamma} \right| + \log (1 + h)
\]

并且 \(\log(1+h)\) 可以展开，如上所述，因为函数 \(\log(1+t)\) 在 0 的一个邻域上是无穷次可导的；此外，若 \(\log g_{\gamma}\) 相对于 E 容许一个渐近展开，或相对于某个尺度 \(E_{1} \supset E\) 容许一个渐近展开，则在把这两个渐近展开相加后得到 \(\log|f|\) 的一个渐近展开。

例. 我们有 \((1+x)^{1/x}=\exp\left(\frac{1}{x}\log(1+x)\right)\) ；当 x 趋于 \(+\infty\) 时，我们有 \(\log(1+x)=\log x+\log\left(1+\frac{1}{x}\right)\) ，由此得 \(\frac{1}{x}\log(1+x)\) 的相对于 \(x^{\alpha}(\log x)^{\beta}\) 的尺度的渐近展开：

\[
\frac {1}{x} \log (1 + x) = \frac {\log x}{x} + \frac {1}{x ^ {2}} - \frac {1}{2 x ^ {3}} + o _ {1} \left(\frac {1}{x ^ {3}}\right).
\]

从这个展开，以及泰勒展开

\[
e ^ {u} = 1 + u + \frac {u ^ {2}}{2} + \frac {u ^ {3}}{6} + o (u ^ {3})
\]

（在 u = 0 的一个邻域上），可导出渐近展开

\[
\begin{array}{r l} (1 + x) ^ {1 / x} & = 1 + \frac {\log x}{x} + \frac {1}{2} \frac {(\log x) ^ {2}}{x ^ {2}} + \frac {1}{x ^ {2}} \\ & \quad + \frac {1}{6} \frac {(\log x) ^ {3}}{x ^ {3}} + \frac {\log x}{x ^ {3}} - \frac {1}{2 x ^ {3}} + o _ {2} \left(\frac {1}{x ^ {3}}\right) \end{array}
\]

（相对于尺度 \(x^{\alpha}(\log x)^{\beta}\) ），用的是上面解释的方法。

保持同样的假设与记号， \(e^{f}\) 的渐近展开除 \(f \gg 1\) 的情形外不产生新问题；这时必须区分两种情形，视 \(g_{\alpha} \gg 1\) 还是 \(g_{\alpha} \preccurlyeq 1\) 而定。在第一种情形，给出 f 的一个展开并不能得到 \(e^{f}\) 相对于 E 的主部，因为一般不知道余项 \(r_{\alpha}\) 是否趋于 0，即 \(e^{r_{\alpha}}\) 是否趋于 1。相反，若 \(g_{\alpha} \preccurlyeq 1\) ，则 \(r_{\alpha} \ll 1\) ，从而 \(e^{f} \sim \exp\left(\sum_{\lambda \leqslant \alpha} a_{\lambda} g_{\lambda}\right)\) 。可以把这个结果弄得更精确：设 \(a_{\gamma} g_{\gamma}\) 是 f 的主部，并设 \(\delta\) 是指标（满足 \(\gamma < \delta \leqslant \alpha\) ，使 \(g_{\delta} = 1\) 者）；我们置 \(f_{1} = \sum_{\lambda \leqslant \delta} a_{\lambda} g_{\lambda}\) ， \(f_{2} = \sum_{\delta < \lambda \leqslant \alpha} a_{\lambda} g_{\lambda} + r_{\alpha}\) ；我们有 \(f = f_{1} + f_{2}\) ，故 \(e^{f} = e^{f_{1}} e^{f_{2}}\) ，而本小节开头解释的一般方法使我们能构成 \(e^{f_{2}}\) 的一个渐近展开（从 \(e^{t}\) 在点 t = 0 处的泰勒展开出发）。于是将得到 \(e^{f}\) 的一个渐近展开，如果 \(e^{f_{1}} = \prod_{\lambda \leqslant \delta} \exp(a_{\lambda} g_{\lambda})\) 属于 E，或属于包含 E 的某个尺度 \(E_{1}\) 。

例. 我们有 \(x^{x^{1/x}} = \exp\left(\log x. \exp\left(\frac{1}{x} \log x\right)\right)\) ；当 x 趋于 \(+\infty\) 时，有 \(\log x \ll x\) ，由此得 \(\log x. \exp\left(\frac{1}{x} \log x\right)\) 的相对于 \(x^{\alpha}(\log x)^{\beta}\) 的尺度的渐近展开：

\[
\log x. \exp \left(\frac {1}{x} \log x\right) = \log x + \frac {(\log x) ^ {2}}{x} + \frac {1}{2} \frac {(\log x) ^ {3}}{x ^ {2}} + o \left(\frac {(\log x) ^ {3}}{x ^ {2}}\right).
\]

这个展开从第二项起的所有项都趋于 0；由这个展开以及泰勒展开 \(e^u = 1 + u + u^2 / 2 + o(u^2)\) （在 \(u = 0\) 的一个邻域上）可导出

\[
x ^ {x ^ {1 / x}} = x + (\log x) ^ {2} + \frac {1}{2} \frac {(\log x) ^ {4}}{x} + \frac {1}{2} \frac {(\log x) ^ {3}}{x} + o \left(\frac {(\log x) ^ {3}}{x}\right).
\]

